% sec1.tex -- writer W1.
\section{Introduction}\label{sec:intro}

\paragraph{The problem.} Identify $\R^2$ with $\C$. A \emph{pursuer path} is a map $H:[0,\infty)\to\R^2$ with $H(0)=0$ and
$|H'|\le1$ a.e., i.e.\ a $1$-Lipschitz path from the origin. A point $S$, the Snitch, starts at distance $d>0$ from the origin in
an unknown direction $\alpha\in[0,2\pi)$ and always moves with speed $v\in(0,1)$ directly away from the pursuer:
\begin{equation}\label{eq:snitch}
S_\alpha'(t)=v\,\frac{S_\alpha(t)-H(t)}{|S_\alpha(t)-H(t)|},\qquad S_\alpha(0)=d\,e^{i\alpha}.
\end{equation}
The pursuer knows $d$ and $v$ but not $\alpha$, and sees the Snitch only when it is within distance $\eps\in[0,d)$ (the
\emph{vision radius}); his path is therefore fixed in advance, i.e.\ it may not depend on $\alpha$. For a continuous $H$ the
right-hand side of \eqref{eq:snitch} is continuous in $t$ and locally Lipschitz in $S$ away from $H(t)$; we call $S_\alpha$, the
\emph{candidate} $\alpha$, the maximal solution, defined up to the first time at which it meets $H$ (since $|S_\alpha'|=v$ and
$\bigl||S_\alpha-H|'\bigr|\le1+v$, a maximal solution defined only on a bounded interval has a limit at its end, where it meets
$H$). Candidate $\alpha$ is \emph{caught} if $S_\alpha(t)=H(t)$ for some $t$ and \emph{seen} if $|S_\alpha(t)-H(t)|\le\eps$ for
some $t$ in its interval of existence or at its end. With $\eps>0$ the pursuer wins once the Snitch is seen; with $\eps=0$ he
must catch it. Since \eqref{eq:snitch} is invariant under $(S,H,t)\mapsto(\lambda S,\lambda H,\lambda t)$, one may take $d=1$.

The question is whether, for given $v$, a single path sees (or catches) \emph{every} direction, and in what time. A path that
runs straight to the circle of radius $d$ and once around it sees every direction when $(1+2\pi)vd\le\eps$: it passes $de^{i\alpha}$ before time $(1+2\pi)d$, when the Snitch is still within $(1+2\pi)vd$ of that point; the difficulty is to obtain a range of
$v$, and a time, that do not degenerate as $\eps\to0$.

\begin{theorem}\label{thm:main}
Let $\vz:=5.8\cdot10^{-3}$, $T_0:=13.1$ and $\epsS:=4\cdot10^{-11}$. For every $d>0$, every $\eps\in(0,d)$ and every
$v\in(0,\vz]$ the path $H=H_{v,\eps,d}$ --- design~T \textup{(Definition~\ref{def:T})} followed by the greedy chase
\textup{(Definition~\ref{def:chase})} on the working front of Theorem~\ref{thm:touchdown}, both for the vision radius
$\eps':=\min(\eps,\epsS d)$ and scaled to distance $d$ --- satisfies $H(0)=0$, $|H'|\le1$ a.e., depends only on $(v,\eps,d)$, and
for every $\alpha\in[0,2\pi)$ there is $t\le T_0d$ with $|S_\alpha(t)-H(t)|\le\eps$. Neither $\vz$ nor $T_0$ depends on $\eps$.
\end{theorem}

\begin{corollary}\label{cor:exact}
Let $d>0$ and $v\in(0,5.8\cdot10^{-3}]$. There is a path $H$ with $H(0)=0$, $|H'|\le1$ a.e., such that for every
$\alpha\in[0,2\pi)$ there is $t\le13.1\,d$ with $S_\alpha(t)=H(t)$.
\end{corollary}

Both are proved in Section~\ref{sec:proof}; the corollary follows from the theorem and the compactness Lemma~\ref{lem:compact},
because the time bound of Theorem~\ref{thm:main} is uniform in $\eps$.

\begin{remark}\label{rem:numerics}
The parameters of Appendix~\ref{app:const} were found by numerical optimisation; the script \texttt{verify\_constants.py}
checks every side condition in $50$-digit arithmetic; the near-pair constant $D'$ (Definition~\ref{def:Dprime}), a supremum
over a two-parameter family, is bounded there by an evaluation in outward-rounded arithmetic. At the chosen parameters several
side conditions are nearly tight at once (Table~\ref{tab:side}), so the method as it stands does not reach much beyond the
stated $\vz$. Simulations suggest that the same path works for much larger $v$.
\end{remark}

\paragraph{Outline of the proof.} We take $d=1$. Inside the vision disc the field $x/|x|$ of \eqref{eq:snitch} is singular, so
we replace it by a $C^1$ field $F$ that coincides with $x/|x|$ outside the disc of radius $\eps$ (Definition~\ref{def:field}).
Curves of \emph{virtual} candidates, moved by $vF(x-P)$, form $C^1$ fronts whose arclength never decreases and whose chords
never shrink (Lemmas~\ref{lem:flow}, \ref{lem:sep}), and a virtual candidate that comes within $\eps$ of the path is a true
candidate that has been seen (Lemma~\ref{lem:seen}). The pursuer can therefore follow such a front without knowing which of
its points carry the Snitch: the \emph{greedy chase} (Definition~\ref{def:chase}) walks along it in straight legs, each aimed
at the point an arclength $\eps$ further on, and every point passed is seen.

The chase finishes in a time independent of $\eps$ only if the front stays nearly straight on every scale between $\eps$ and a
fixed length $R$. This is measured by the two-point \emph{excess} $E=l-c$, arc length minus chord, of pairs of labels
(Definition~\ref{def:excess}); its evolution under the flow is governed by the exact identities of Lemma~\ref{lem:excess}, in
which the stretching $T\cdot AT$ of the front and the turning of its chord appear explicitly. The end theorem
(Theorem~\ref{thm:end}) is a bootstrap: if at the start the excess of a pair of label distance $\lambda$ is at most a small
multiple of $\lambda^{1+2\gamma}$, with $\gamma>\frac12$, and chords are comparable to $\lambda$, then the length bound
$l\le\kappa\lambda$ and the excess bound $N=(2E)^{1/2}\le M\lambda^{1/2+\gamma}$ persist along the whole chase, every
leg advances the label by at least $\eps/\kappa$, and the chase ends after time $\kappa U/(1-v)$, where $U$ is the length of
the front. Pairs far ahead of the pursuer are controlled by a Duhamel inequality (Lemma~\ref{lem:duhamel}), pairs near him by
integrating the excess identity over the legs, with the pursuer's position averaged along each leg (Section~\ref{sec:farnear}).
No constant depends on $\eps$.

The start data come from \emph{design~T} (Definition~\ref{def:T}). The pursuer first runs straight through the set of possible
positions, catching direction $0$ and seeing a thin sector around it; he then turns and runs along a straight line that
touches the remaining front tangentially, at the label $\alpha_0$, at the touch-down time $\ttd$. Along a straight run the
motion of the candidates has two exact integrals (Proposition~\ref{prop:runs}), and their derivatives in the direction label
transport the tangent and the curvature of the front exactly (Lemma~\ref{lem:transport}). This yields the front at touch-down
explicitly enough (Theorem~\ref{thm:touchdown}): it is a $C^1$ curve of bounded length that starts at the pursuer, is
flat at \emph{both} ends with tangent close to the run-2 heading, has excess $N\le c_N\lambda^{3/2-\muTb}$ on small scales and
chords bounded below on large ones (Theorem~\ref{thm:touchdown}, Lemma~\ref{lem:T2star}).

The front ends close to its start: its last point lies a distance at most of order $\eps^{1/2}$ behind the pursuer at $\ttd$. Pairs
with one point at the beginning of the front and the other near its end do not satisfy the start conditions. They are
\emph{excused}: the end theorem claims nothing about them until the chase has advanced by a fixed length $\We$. Meanwhile the
pursuer moves away from the end of the front with a controlled heading (Lemma~\ref{lem:heading}, obtained from a second
application of the end theorem to a short initial arc), and the departure theorem (Theorem~\ref{thm:departure}) shows that the
excused pairs are then far pairs with the start data the bootstrap needs. Assembling these pieces with explicit constants
(Section~\ref{sec:proof}, Appendix~\ref{app:const}) gives Theorem~\ref{thm:main}; letting $\eps\to0$ along a sequence of these
paths gives Corollary~\ref{cor:exact}.

\paragraph{Organisation.} Section~\ref{sec:prelim} collects the calculus of virtual fronts and of the excess. Section~\ref{sec:end}
defines the chase and states and organises the proof of the end theorem, which is completed in Section~\ref{sec:farnear}.
Section~\ref{sec:dep} proves the heading lemma and the departure theorem. Section~\ref{sec:designT} treats straight runs and
defines design~T, and Section~\ref{sec:touch} analyses the front at touch-down. Section~\ref{sec:proof} assembles the proof, and
Appendix~\ref{app:const} lists the constants and the side conditions they satisfy.

\paragraph{Notation.} Vectors in $\R^2=\C$: $a\cdot b:=\operatorname{Re}(\bar ab)$, $a\times b:=\operatorname{Im}(\bar ab)$,
$a^\perp:=ia$; thus $a^\perp\cdot b=a\times b$. Throughout Sections~\ref{sec:prelim}--\ref{sec:proof} (except in the last two proofs
of Section~\ref{sec:proof}) $d=1$, $0<v\le\vz$, $0<\eps\le\epsS$ (in Sections~\ref{sec:designT} and~\ref{sec:touch}:
$0<\eps\le10^{-3}$), and $\mu:=v/(1-v)$. Named constants are numbers, listed in Appendix~\ref{app:const}. The notation for fronts is
fixed in Section~\ref{sec:prelim}.
